% !TEX program = lualatex
% Chapter 5: Constitutional Safeguards Audit, Anti-Censorship Review & Democratic Governance
% CivicShield AI Preliminary Report — Government-Tender Edition

\chapter{ביקורת ערובות חוקתיות, מניעת צנזורה וממשל דמוקרטי}
\label{ch:safeguards_audit}

\section{מבוא ומסגרת הביקורת החוקתית}
\label{sec:constitutional_framework}

הפעלתה של מערכת מבוססת בינה מלאכותית במרחב הציבורי והדיגיטלי להגנה מפני מבצעי השפעה זרים, רשתות בוטים מתואמות ודיסאינפורמציה ממומנת מחייבת עמידה ברף החוקתי והמנהלי הגבוה ביותר. מערכת \textenglish{CivicShield AI} פועלת בצומת רגיש שבין הגנה על ריבונות המידע והתהליך הדמוקרטי לבין שמירה קפדנית על זכויות יסוד, ובראשן חופש הביטוי, הזכות לפרטיות, חופש ההתאגדות והשוויון בפני החוק.

ביקורת הערובות החוקתיות בפרק זה נשענת על עיקרון מרכזי אחד: \textbf{הגנה על הדמוקרטיה אינה יכולה להישען על הבטחות מדיניות בלבד, אלא חייבת להיות מוטמעת בארכיטקטורה הטכנולוגית ובמנגנוני הממשל} (\textenglish{Constitutional-by-Design}). בהתאם לעקרונות המשפט החוקתי הישראלי, ובראשם חוק יסוד: כבוד האדם וחירותו~\cite{basiclaw1992}, הלכת קול העם~\cite{kolhaam1953} ומבחני פסקת ההגבלה (הסמכה מפורשת בחוק, תכלית ראויה ומידתיות על שלושת מבחני המשנה שלה~\cite{barak2012proportionality}), כל רכיב במערכת נבחן מול שלוש שאלות יסוד:
\begin{enumerate}
    \item \textbf{סמכות ותיחום:} האם המערכת מוגבלת לניתוח תבניות התנהגות מתואמות וגלויות (\textenglish{OSINT}) בלבד, ללא חריגה לאיסוף מידע פרטי או לאכיפה עצמאית?
    \item \textbf{מידתיות וצמצום פגיעה:} האם קיים מנגנון טכני המבטיח כי הפגיעה האפשרית בשיח ציבורי לגיטימי שואפת לאפס, וכי כל התרעה מדורגת לפי רמת ודאות וסיכון?
    \item \textbf{ביקורתיות והפיכות:} האם כל תוצר של המערכת ניתן לשחזור, להסבר (\textenglish{Explainability}), לביקורת חיצונית ולערעור?
\end{enumerate}

\section{ביקורת מניעת צנזורה והגנה על שיח פוליטי לגיטימי}
\label{sec:anti_censorship}

הסיכון החוקתי החמור ביותר במערכות לניטור מרחב מידע הוא גלישה מהגנה מפני התערבות זרה ומתואמת אל עבר תיוג שגוי או דיכוי של שיח אזרחי אותנטי, ביקורת פוליטית, סאטירה או מחאה חברתית. כדי למנוע תרחיש זה באופן מבני, אומצו במערכת \textenglish{CivicShield AI} ארבעה קווי הגנה מחייבים נגד צנזורה:

\subsection{עקרון ההפרדה המבנית: התנהגות מתואמת מול תוכן רעיוני}
המערכת \textbf{אינה} מוסמכת ואינה מתוכננת להכריע בשאלת "האמת" של טענות פוליטיות או אידיאולוגיות. מנוע הזיהוי מתמקד בחתימות התנהגותיות ורשתיות של התנהגות בלתי-אותנטית מתואמת (\textenglish{Coordinated Inauthentic Behavior -- CIB}), כגון סנכרון זמני חריג בין מאות חשבונות, הפצה מתוזמרת ממקורות תשתית משותפים, דפוסי יצירת חשבונות סינתטיים ומניפולציה אלגוריתמית של חשיפה. תוכן ההודעה משמש לכל היותר לצורך אשכול נושאי והערכת הקשר, ולעולם אינו משמש עילה עצמאית לסימון חשבון או פרסום.

\subsection{אי-וריאנט אפס צנזורה ואפס הסרה אוטומטית}
למערכת \textenglish{CivicShield AI} \textbf{אין} הרשאת כתיבה, מחיקה, חסימה או הגבלת תפוצה באף פלטפורמה חברתית או ערוץ תקשורת. המערכת פועלת במודל \textenglish{Read-Only Advisory}: תוצריה הם התרעות מודיעיניות, דוחות תמונת מצב ומדדי סיכון המוגשים לגורמים אנושיים מוסמכים בלבד. כל פעולה אופרטיבית מול פלטפורמות או מול הציבור (כגון פרסום הבהרה רשמית או פנייה משפטית) מחייבת אישור אנושי מתועד (\textenglish{Human-in-the-Loop -- HITL}) וחוות דעת משפטית בהתאם לדין ולעקרונות פסיקת בג"ץ בעניין יחידת הסייבר~\cite{adalah2021cyber}.

\subsection{מיפוי שגיאות חיוביות כוזבות (\textenglish{False Positives}) בשיח אזרחי}
במסגרת ביקורת הצוות האדום (\textenglish{Red-Team}) זוהו ארבע קטגוריות רגישות במיוחד לשגיאות סיווג, שלגביהן הוגדרו חסמי סף מחמירים ומבחני רגרסיה קבועים:
\begin{itemize}
    \item \textbf{מחאה אזרחית וקמפיינים פוליטיים לגיטימיים:} שיתוף המוני של סיסמאות או האשטאגים זהים בתוך פרק זמן קצר עשוי להידמות סטטיסטית לתיאום בוטים. המערכת נדרשת להבחין בין תיאום אורגני גלוי (כגון קמפיין מפלגתי או ארגון חברה אזרחית מזוהה) לבין תיאום סמוי ומזויף.
    \item \textbf{סאטירה, פרודיה ושיח בוטה:} שפה חריפה, צינית או מתריסה מוגנת תחת חופש הביטוי הפוליטי. מודלי השפה נבחנים באופן ייעודי כדי למנוע תיוג של סאטירה כדיסאינפורמציה.
    \item \textbf{קבוצות מיעוט ורב-לשוניות:} בישראל מתקיים שיח ציבורי ער בעברית, ערבית, רוסית, אמהרית ואנגלית. ביקורת הצוות האדום מחייבת שוויון בשיעורי השגיאה (\textenglish{Equalized False Positive Rate}) בין שפות ודיאלקטים, כדי למנוע פגיעה בלתי-מידתית בקהילות מיעוט.
    \item \textbf{עיתונות חוקרת ומדליפים:} הפצה ראשונית של מידע רגיש על ידי עיתונאים או אזרחים אינה מסווגת כחריגה עוינת, ומנגנון ההגנה כולל רשימות החרגה שקופות ונוהל בדיקה מואץ.
\end{itemize}

\subsection{ניטרליות פוליטית וביקורת הטיה השקפתית (\textenglish{Viewpoint Neutrality})}
כל גרסת מודל עוברת מבחן סימטריה אידיאולוגית לפני פריסה: אותם דפוסי תיאום והפצה נבחנים על גבי קורפוסים סינתטיים המייצגים עמדות מנוגדות לאורך הקשת הפוליטית. פער העולה על \textenglish{2\%} בשיעור ההתרעות בין עמדות סימטריות פוסל את הגרסה אוטומטית ומונע את עלייתה לסביבת הייצור.

\section{עקרונות הפעלה חוקתיים ומגבלות סמכות}
\label{sec:constitutional_constraints}

טבלה~\ref{tab:safeguards_matrix} מרכזת את מפת הסיכונים החוקתיים של המערכת, הזכויות המוגנות, הערובות הטכנולוגיות והמנהליות שהוטמעו, וכן את מדדי האימות המחייבים לצורך קבלה במכרז הממשלתי.

\begin{table}[htbp]
\centering
\small
\caption{מטריצת ערובות חוקתיות, מגבלות סמכות ומדדי אימות למערכת \textenglish{CivicShield AI}}
\label{tab:safeguards_matrix}
\renewcommand{\arraystretch}{1.25}
\begin{tabularx}{\textwidth}{p{2.7cm} p{2.5cm} X p{3.2cm}}
\toprule
\rowcolor{CSNavy}
\thead{סיכון חוקתי} & \thead{זכות יסוד מושפעת} & \thead{ערובה טכנולוגית ומבנית} & \thead{מדד אימות וסף קבלה} \\
\midrule
צנזורה עקיפה או אפקט מצנן על שיח פוליטי & חופש הביטוי הפוליטי~\cite{kolhaam1953} & ארכיטקטורת \textenglish{Read-Only}; איסור הסרה אוטומטית; הפרדת \textenglish{CIB} מתוכן & \textenglish{0} פעולות חסימה אוטומטיות; פערי הטיה השקפתית $< 2\%$ \\
פרופיילינג ומעקב אחר אזרחים פרטיים & הזכות לפרטיות וכבוד האדם~\cite{privacy_law1981} & אנונימיזציה במקור; איסור בניית תיק אישי לאזרח; עיבוד מידע פומבי בלבד & \textenglish{100\%} תאימות למדיניות צמצום מידע; \textenglish{TTL} מחיקה $\le 30$ יום לנתונים גולמיים \\
אפליה שפתית או מגזרית בזיהוי & שוויון ואי-אפליה & מבחני הוגנות רב-לשוניים (עברית, ערבית, רוסית, אנגלית) בכל עדכון מודל & פער \textenglish{FPR} בין שפות וקהילות $< 3\%$ \\
שימוש לרעה במערכת לצרכים פוליטיים & טוהר המידות וניטרליות השירות הציבורי & בקרת גישה מבוססת מאפיינים ותפקידים (\textenglish{ABAC/RBAC}), אישור כפול (\textenglish{Two-Person Rule}) ויומן בלתי-ניתן לשינוי & \textenglish{100\%} שאילתות מתועדות בחתימה קריפטוגרפית (\textenglish{Hash-Chained Audit Log}) \\
החלטות "קופסה שחורה" ללא יכולת ערעור & הגינות מנהלית וזכות הטיעון~\cite{eu_aiact2024} & הפקת דוח הסבר מובנה לכל התרעה; שמירת ראיות תשתית וגרסת מודל & \textenglish{100\%} התרעות כוללות תעודת מקור (\textenglish{Provenance}) ורמת ודאות \\
\bottomrule
\end{tabularx}
\end{table}

\section{אחריותיות דמוקרטית, פיקוח פרלמנטרי ומנגנוני ממשל}
\label{sec:democratic_governance}

לגיטימציה ציבורית היא תנאי הכרחי להצלחת מערכת הגנה על מרחב המידע. בהיעדר אמון ציבורי, גם התרעה מדויקת מפני התערבות זרה עלולה להיתפס ככלי פוליטי. לפיכך, מוצע מודל ממשל ואחריותיות רב-שכבתי (\textenglish{Multi-Tier Democratic Governance}) המבוסס על ארבע שכבות פיקוח בלתי-תלויות:

\subsection{שכבה א': בקרה פנימית וצוות אדום עצמאי (\textenglish{Internal Red-Team \& Ethics Board})}
במסגרת הגוף המפעיל יפעל צוות אדום פנימי הכפוף ישירות ליועץ המשפטי ולקצין הגנת הפרטיות, ולא למנהל המוצר או לראש צוות הפיתוח. לצוות זה סמכות להשהות פריסת מודלים (\textenglish{Deployment Veto}) במקרה של חריגה ממדדי ההטיה, הפרטיות או העמידות.

\subsection{שכבה ב': ועדת פיקוח ציבורית בלתי-תלויה}
תוקם ועדת פיקוח ציבורית בראשות שופט בדימוס, ובה חברים נציגי אקדמיה בתחומי המשפט החוקתי ואתיקה של בינה מלאכותית, נציגי ארגוני זכויות אדם וחופש ביטוי, ומומחי אבטחת מידע. לוועדה תינתן גישה מלאה ליומני הביקורת, למדדי השגיאה, ולמדגם אקראי של התרעות שהופקו, תוך שמירה על סיווג ביטחוני מתאים היכן שנדרש.

\subsection{שכבה ג': פיקוח פרלמנטרי וביקורת מבקר המדינה}
המערכת תפעל תחת דיווח תקופתי (רבעוני ושנתי) לוועדה ייעודית של הכנסת (ועדת המשנה למודיעין או ועדת החוקה, חוק ומשפט) ותעמוד לביקורת שוטפת של מבקר המדינה והרשות להגנת הפרטיות. שינוי במנדט המערכת, במקורות המידע המאושרים או בהגדרות הסף המרכזיות יחייב עדכון פורמלי ואישור של ועדת הפיקוח.

\subsection{שכבה ד': שקיפות ציבורית ומנגנון ערעור ותיקון (\textenglish{Transparency \& Redress})}
אחת לחצי שנה יפורסם לציבור \textbf{דוח שקיפות פומבי} שיכלול: מספר ההתרעות שנוצרו לפי קטגוריה, שיעורי שגיאה חיובית כוזבת שהתגלו בדיעבד, מספר הפניות שהועברו לפלטפורמות ותוצאותיהן, וסיכום ממצאי ביקורת הצוות האדום. בנוסף, כל גוף תקשורת, ארגון אזרחי או אזרח שסבור כי תוכן או פעילות שלו תויגו בשגגה יוכל להגיש בקשת בירור וערעור לנציב תלונות עצמאי, המחויב למענה מנומק בתוך \textenglish{14} ימי עבודה ולתיקון רטרואקטיבי של המערכת במקרה של טעות.

\section{מתודולוגיית צוות אדום (\textenglish{Red-Team}) ובדיקות עמידות}
\label{sec:redteam_methodology}

ביקורת הערובות החוקתיות אינה אירוע חד-פעמי בשלב ההקמה, אלא תהליך מתמשך לאורך מחזור החיים של המערכת. מתודולוגיית ה-\textenglish{Red-Team} של \textenglish{CivicShield AI} פועלת במחזורים רבעוניים וכוללת ארבעה צירי תקיפה ובדיקה:
\begin{enumerate}
    \item \textbf{מבחני הטיה וצנזורה יתר (\textenglish{Over-Triggering Stress Tests}):} הזרקת תרחישי שיח אזרחי סוער, מחאות המוניות, סאטירה וביקורת חריפה נגד מוסדות שלטון, כדי לוודא שהמערכת אינה מפיקה התרעות שווא.
    \item \textbf{מבחני התחמקות והסתגלות יריב (\textenglish{Adversarial Evasion}):} הדמיית מבצעי השפעה מתקדמים המשתמשים במודלי שפה גדולים (\textenglish{LLMs}), התחזות סגנונית ופיזור זמנים, כדי לבחון את חוסן המערכת מבלי להוריד את רף ההגנה על שיח לגיטימי.
    \item \textbf{מבחני הרעלת נתונים וסחיפת מודל (\textenglish{Data Poisoning \& Model Drift}):} ניטור מתמשך של התפלגות ההתרעות לאורך זמן, במיוחד בתקופות רגישות כמו מערכות בחירות או מצבי חירום, לזיהוי סחיפה סטטיסטית או ניסיון מכוון לשבש את הסיווג.
    \item \textbf{נוהל מפסק חירום (\textenglish{Circuit-Breaker / Kill-Switch}):} הגדרת תנאי סף אוטומטיים להשעיית מודול סיווג (למשל, קפיצה חריגה של מעל \textenglish{300\%} בשיעור ההתרעות בתוך \textenglish{24} שעות ללא אימות אנושי) והעברת המערכת למצב בדיקה מוגן.
\end{enumerate}

\section{סיכום ממצאי הביקורת והתחייבויות למכרז}
\label{sec:safeguards_summary}

ביקורת הערובות החוקתיות מעלה כי ארכיטקטורת \textenglish{CivicShield AI} עומדת בדרישות הסף החוקתיות והדמוקרטיות של מכרז ממשלתי, בתנאי ששלוש התחייבויות ליבה יישמרו כדרישות חובה בלתי-ניתנות לוויתור (\textenglish{Non-Negotiable Invariants}) בכל שלבי הפיתוח והתפעול:
\begin{itemize}
    \item \textbf{איסור מוחלט על אכיפה אוטומטית:} המערכת תישאר כלי התרעה וניתוח בלבד (\textenglish{Advisory-Only}) עם שיקול דעת אנושי ומשפטי מחייב.
    \item \textbf{שקיפות ובקרה חיצונית כתנאי להפעלה:} פריסה מבצעית מותנית בהקמת ועדת הפיקוח הציבורית, הפעלת יומני ביקורת חתומים קריפטוגרפית ופרסום דוח שקיפות חצי-שנתי.
    \item \textbf{עמידה מתמשכת במדדי אי-הטיה:} כל עדכון תוכנה או מודל יחויב במעבר מוצלח של סדרת מבחני ה-\textenglish{Red-Team} הרב-לשונית והסימטרית-פוליטית כמפורט בפרק זה.
\end{itemize}
